\chapter*{Prefazio}
\addcontentsline{toc}{chapter}{Prefazio}

Questo lavoro nasce dall’esigenza di restituire il testo del \textit{Śrī Lalitā Sahasranāma} nella sua densità originaria, senza riduzioni devozionali né sovrapposizioni interpretative estranee. Ogni śloka è presentato secondo una struttura costante: testo sanscrito attestato, traslitterazione, anvaya, resa italiana e commento essenziale. L’intento non è quello di esaurire il significato del testo, ma di renderlo leggibile nella sua coerenza interna, lasciando che sia la sequenza stessa dei nāma a guidare la comprensione.

L’approccio adottato è volutamente sobrio. I commenti non introducono simbolismi aggiuntivi né sistemi dottrinali non richiesti dal testo, ma si limitano a chiarire la funzione dei nomi e delle immagini così come emergono dalla loro posizione e dal loro rapporto reciproco. In particolare, si è prestata attenzione a distinguere i nāma effettivi dalle qualificazioni descrittive, e a mantenere una progressione che rispetti il passaggio dalla sovranità ontologica all’azione cosmica, fino alla manifestazione sensibile della forma.

La lingua italiana utilizzata mira alla precisione più che all’eleganza, privilegiando la chiarezza concettuale e la continuità terminologica. I termini sanscriti sono mantenuti quando la traduzione rischierebbe di appiattirne il valore semantico, e segnalati tipograficamente in modo coerente lungo tutto il testo.

Questo volume non pretende di sostituire la tradizione commentariale, né di offrire una lettura definitiva. Si propone piuttosto come strumento di studio e di contemplazione ordinata, rivolto a chi desidera accostarsi al \textit{Sahasranāma} con rispetto per il testo e attenzione alla sua architettura interna.
